% !TeX root = ../papers/nc-p3.tex
\section{Bell曲面の全単純スペクトルと可逆二平面分解の分類}
\label{sec:split-simple-surface}
\begin{theorem}\label{thm:split-simple-surface}
$x=(1,-uv,u,v)$, $f(t)=(t^2-1)/(2t)$, $A=f(u)^2$, $B=f(v)^2$ とする。
平方が相異なる可逆曲面部の基底は $uv\ne0$, $AB\ne0$, $A\ne B$ である。
以下が全ての無印添字保持同型類である。
\begin{enumerate}
\item 通常の四表示を $h=1,X,Z,XZ$ の順に $0,1,2,3$ と標識する。
$(u,v,h)$ に対して以下の三変換とそれらの合成が全ての同型を与える。
右欄は旧番号 $0,1,2,3$ に対する新番号の順である。
\[
\begin{array}{c|c}
(u,v)\mapsto(-u,v)&3,2,1,0\\
(u,v)\mapsto(1/u,v)&1,0,3,2\\
(u,v)\mapsto(v,u)&2,1,0,3
\end{array}
\]
各無順序対 $\{A,B\}$ の上には、どちらも $-1$ と異なれば四類、
一つが $-1$ なら二類がある。全て最小周期2で、十五次元成分に属する。
\item $\{A,B\}=\{-1,b\}$, $b\ne0,-1$ のとき、各 $b$ にさらに一類がある。
これは十五次元成分に属さず、最小周期は $b=1$ で1、他で2である。
この一変数族の基底軌道の閉包は、制約なしの階数スキームの射影60次元既約成分で、
各 $b\ne0,-1$ の軌道で滑らかである。
\end{enumerate}
いずれも標準Koszul左右AS正則かつ多重安定である。
\end{theorem}
\begin{proof}[基底の同型と移送商]
平方の集合は $\{1,u^2,v^2,u^2v^2\}$ であり、可逆性の下では平方の衝突は $A=B$ と同値である。
この点は三つのBell曲面のうちただ一つにある。例えば $ab+cd=ac+bd=0$ なら
非零座標から $b^2=c^2$ を得る。ゆえに曲面内の基底比較は
$ab+cd=0$ を保つBell付け替えに限られる。四座標の二対分割を保つ八置換と射影偶符号四通りで位数32、
$(u,v)$ 上では独立な符号・逆数と交換である。その作用の生成は下で実際の二文字基底でも与える。
従って基底の完全不変量は $\{A,B\}$ である。

左右の直交四次元模型で単純平方スペクトルを持つため $H$ のうちこれらの固有線を個別に保つ部分は
同時の二文字Pauli四元群 $H_0$ である。固有線に非自明な置換が生じるなら
$P$ の射影倍率の四乗は1で、その二乗は $-1$ でなければならない。
平方集合が負号で不変となるのは $u^2=-1$ または $v^2=-1$ のときだけである。
実際 $u^2v^2=-1$ だけで他の二平方も負号で組を作れば $u^4=1$ となる。
従って一般には $H=H_0$、通常の四対応 $(hZ,h,hZ,h)$ が四共役類を尽くす。

例外では $u=\iota$, $\iota^2=-1$ に取れる。$E=(1,X,Z,XZ)$ が $P$ の射影直線を保ち、
$H$ は $H_0$ と $E$ の生成する位数8の可換群である。位数の上界は固有線を個別に保つ核と
可能なただ一つの負号置換による。$D_0=(Z,1,Z,1)$ とすると通常基底の第二対応全体は $H D_0$ である。
$H_0$ は比較作用で自明に働くが、$E$ は四箇所同じ $Z$ を掛ける。
従って $D_0,FD_0,ED_0,FED_0$ が印付きの相異なる四類である（$F$ は全位置の $X$）。
分解交換では第一交代成分が逆数Bell四数に移る。
$(1,Z,Z,1)$ でこれを元に戻すと、通常表示は各々自分に戻る
（定理\ref{thm:split-surface-full-fibers}の基底等式は特殊点でも成立する）。
追加表示は $ED_0\leftrightarrow FED_0$ となる。
$H$ は有限で分解トーラスが階数1の極大トーラスなので、分解の符号と共役を除いて追加同型はない。

通常表示において、三つの基底生成変換を持ち上げる能動基底を順に
$(1,1,XZ,XZ),(1,X,X,1),(D,D^{-1},D^{-1},D)$, $D=\operatorname{diag}(1,\iota)$ とする。
各々を $P$ と四対応に作用させると表を得る。
\path{computations/verify_split_simple_surface.py} はこの三つの射影等式を $\mathbb Q(\iota,u,v)$ で検算する。
付け替え群の作用は一般の $(u,v)$ で自由であり、$u^2=-1$ の固定変換 $u\mapsto-1/u$ は
ちょうど同時 $Z$ による二番号の同一視を与える。これで無印分類を得る。
各 $P,Q$ は補語に等しい係数を持つBell型、四対応はPauli型なので、
各一文字のGramの非対角成分が消え、各ブロック内の両ビットの周辺は等しい。
既出の六ブロックのノルム最小化により多重安定であり、正則性は可逆分解の定理による。
\end{proof}

\begin{proof}[全ての一移動]
$\phi(t)=(t-1)/(t+1)$ とすると、巡回したBell四数を先頭1に直したパラメータは
$(-1/\phi(v),\phi(u))$ となる。一移動と分解交換では第二対応が $L_{i-1}^{-1}$ となり、
通常表示の番号は同時 $Z$ 倍になる。二移動で第一成分は元に戻り、第二成分は通常表示で一致し、
追加表示でも $E$ の比較で消せる同時 $Z$ 倍にしか変わらない。
一移動があり得るのは $\{A,B\}=\{1/A,1/B\}$ の場合だけである。
$AB=1$ かつ $A\ne B$ なら、付け替えで $v=\phi(u)$ に取れる。
このとき一移動後の $(u,v)$ は既に元通りだが、番号が $Z$ 倍で変わり、基底に固定変換はない。
残りの候補は $\{A,B\}=\{-1,1\}$ だけで、$u=\iota,v=1+\sqrt2$ に取る。
一移動は $(-v,\iota)$ と $Z$ 倍へ移り、交換と第二変数の負号で基底を戻すと、
二つの通常番号類 $\{0,2\}$ と $\{1,3\}$ が入れ替わる。
追加類はその基底上で唯一の追加類なので固定される。これで全周期を得る。
\end{proof}

\begin{proof}[追加表示の成分]
通常表示は定理\ref{thm:split-surface-component}による。追加表示では分解の交換が
印付きの\emph{別の}類へ移るので、この窓に指定分解を逆転する自己同型はない。
射影安定群は多重安定性から簡約的であり、もし恒等成分が階数1のトーラスより大きければ
非自明なWeyl反射がこれを逆転する。従って射影軌道次元は $60-1=59$ である。
$b$ は非定数でその上に一類なので族の閉包は射影60次元である。

$u=\iota$, $v(v^4-1)\ne0$ において $ED_0$ の二切断のアフィン接次元は常に61である。
上の検算では二色内の純関係のビットに合わせ、六降順座標の射影で関係を表示し、
二切断の左右の核の全200積を取る。分母は $v(v^4-1)$ の局所化で単元である。
同一方程式に現れる列をつないでブロックを分け、
$\mathbb Q(\iota)[v]$ 上のSmith形で非零不変因子の根が全て許容域外にあることと、
核次元和61を調べる。これは上の軌道族のアフィン次元であるから全点で滑らかで、
通る成分はこの軌道族の閉包一つだけである。特にアフィン76次元の既知成分の点はない。
\end{proof}

\begin{corollary}\label{cor:all-invertible-split-atlas}
可逆二平面分解を持つ全窓の添字保持同型類と最小周期は、
直線の半単純表・二固有値Jordan表・単一固有値Jordan表、曲面の重複表・単純表、
および六次式の一点で尽きる。全ての最小周期は1または2である。
\end{corollary}
\begin{proof}
補完可能な全スペクトルと鎖分割は定理\ref{thm:split-bell-locus}、
\ref{thm:split-jordan-transport}の三領域で尽くされる。
曲面と直線の共通部分は直線の表で処理した。
異なる印があり得るときは、極大トーラスでの分解判定、または四指標が異なるときの単項同型判定を
各表で使った。量子窓では座標四文字を二対に分ける他の印もあるが、
どの可逆分解も、対応する一径数を量子窓の共通対角極大トーラスへ共役すれば二対分割に戻るので、
その交代成分は常に半単純直線型である。表外の重複は生じない。六次式の一点とは第三スペクトルで区別できる。
\end{proof}
